% Modernised reading — fonds Alexandre Grothendieck, shelfmark 107, pages 1-25
% (the whole folder).
%
% This is an INTERPRETATION, not a transcription. It re-reads the pages in
% today's notation and vocabulary, states the hypotheses the manuscript leaves
% implicit, and drops the critical apparatus so that the mathematics reads
% straight through. Everything the brackets and underlines carried moves into
% a footnote. It is held to being mathematically correct as it stands; where
% the manuscript is loose or mistaken, the true statement is what appears here
% and the footnote says what the page has. In any dispute of substance the
% transcription governs: whoever cites Grothendieck must cite
% batch-01.fr.tex and batch-02.fr.tex.
%
% Scope: the folder was read whole — two batches, pages 1-25 — before any of
% this was written, and it is ONE file for the shelfmark. Pages 4, 6, 8, 10,
% 12, 19 and 25 are printed listings with nothing of his mathematics; pages 1
% and 17 are covers.
%
% Notation fixed for the whole document.
%  - p_A : A x B -> B and p_B : A x B -> A, his convention: the index names
%    the factor the projection forgets. Hom without index (underlined) is the
%    internal Hom of a presheaf category; Hom_A = p_{B*} Hom and
%    Hom_B = p_{A*} Hom are the relative Homs, valued in A^ and B^.
%  - X' : A° -> B^ and X'' : B° -> A^ the two curried forms of X in (A x B)^
%    (page 3's convention, kept on page 5 where he swaps them).
%  - L ⊠ M = p_B^*(L) x p_A^*(M); X^M = Hom(p_A^*(M), X); X ⊗ M the levelwise
%    copower; Gamma''_M(X) = p_{B*} X^M.
%  - pi_0(P) for a presheaf P: the set of connected components of its
%    category of elements (the edition's, used to state the corrected
%    Proposition of page 13).
%  - In the second part: Hot_A, Hot(W) as he writes them; J his interval
%    object, X ⊗ J the cylinder (X x Delta(1) when A = Delta).
%
% Substantive departures from the pages (footnoted where they occur):
%  - p. 13: the Proposition (L, M) |-> L ⊠ M fully faithful, with
%    Hom(L' ⊠ M', L ⊠ M) = Hom(L', L) x Hom(M', M), is false in general
%    (M' = empty); the true formula carries exponents pi_0(M'), pi_0(L'), and
%    the statement holds for L', M' connected. The domain is A^ x B^, not its
%    opposite. The Corollary on internal Homs needs a x L', b x M' connected.
%  - p. 13: the last identification with a single « ! » on (A^ x B^)° would
%    describe A^ x B^, not (A x B)^; the separate condition !! is the right one.
%  - p. 15: p_B^* commutes with internal Homs (true in general, Frobenius);
%    but Hom(L' ⊠ M', L ⊠ e) = Hom(L', L) ⊠ e for arbitrary M', and L^M = L
%    for arbitrary M, are false; true when b x M is connected for every b
%    (for B = Delta: M connected). The canonical map L -> L^M always exists.
%  - p. 9: the comparison map between the coproduct in A^ and in A goes from
%    the former to the latter; the page draws it the other way.
%  - p. 11: ^gX(a) needs p_{A*}, not p_{B*}; the underbrace X^M labels
%    p_{B*} X^M = Gamma''_M(X).
%  - p. 5: Gamma''_M(X) = epsilon_M o X', not epsilon_M o Gamma''_M, and X'
%    (not X'') is the form A° -> B^.
%  - p. 20: « Injectivity of (*) ? » introduces what is the surjectivity of
%    (*); page 24 records it as such.
%  - p. 22: the struck « non » before « trivial cofibration »: j is a
%    cofibration and not a trivial one.
%  - p. 24 and the margin of p. 20: the injectivity of (*) fails in general
%    (Milnor's lim^1; an example is given, the edition's), so X_oo is not a
%    colimit in Hot and A^ -> Hot_A does not commute with countable filtered
%    colimits; the rectification criterion of p. 20's margin is right and can
%    fail (an example is given, the edition's).
%
% Announced and not established: the answer to the question of page 18 (the
% folder leaves it at « not clear at all … Ask Ronnie Brown! »); the
% interpretation of Hom(M, a) => a (p. 14, « ?? »); the sense of page 2's
% sentence that breaks off; Quillen's proposition (Chap. II, n° 2.4) is cited,
% not reproduced, and was not checked against Quillen's text here. No page
% carries a letter c of the English run.
%
% Pass: Opus 5.5 (claude-opus-5-5), 2026-10-03 — first pass, unchecked against the pages by a human.
\documentclass[11pt,a4paper]{article}
\input{../preamble/grothendieck}

\folder{107}
\pages{1}{25}
\dating{[à partir de 1982]}
\watermark{Édition de démonstration\\interprétation personnelle de l'œuvre}
\foldertitle{Closed models [notes postérieures à PS ?] : notes manuscrites (s.d.) — lecture modernisée du dossier entier}

\begin{document}

\section*{Résumé}

\begin{resume}
Ces vingt-cinq feuillets, écrits au dos de listings d'ordinateur, sont des
notes de travail sur l'homotopie vue à travers les « modèles fermés » de
Quillen, c'est-à-dire les catégories de modèles. Ils ont deux parties.

La première monte une machinerie. Une famille d'objets indexée par deux
catégories $A$ et $B$ --- un tableau à double entrée --- peut se lire ligne par
ligne ou colonne par colonne, et l'on a le droit de passer d'une lecture à
l'autre. Le cas qui intéresse Grothendieck est celui où $B$ est la catégorie
des simplexes : un objet du tableau est alors un objet simplicial, une suite
d'objets $X_{0}, X_{1}, X_{2}$, …, reliés par des faces et des
dégénérescences, comme les sommets, arêtes, triangles d'une triangulation. Il
s'agit d'exprimer, dans ce langage, deux opérations : « élever $X$ à la
puissance d'une forme $M$ » --- l'objet des familles de points de $X$
paramétrées par $M$ --- et « multiplier $X$ par $M$ », c'est-à-dire en prendre
autant de copies que $M$ a de cellules, recollées comme $M$. Ces deux
opérations sont adjointes l'une de l'autre, à la manière dont, pour des
ensembles, une application de $X \times M$ vers $Y$ est la même chose qu'une
application de $X$ vers l'ensemble des applications de $M$ dans $Y$. C'est
exactement la structure dont Quillen a besoin pour parler de catégories de
modèles « simpliciales », dans lesquelles les homotopies se manipulent comme en
topologie.

La seconde partie pose une question précise : la catégorie homotopique, où
l'on a identifié les objets ayant le même type d'homotopie, a-t-elle des
limites inductives le long d'une suite $X_{0} \to X_{1} \to X_{2} \to \cdots$,
et la limite ordinaire la fournit-elle ? Une application de la limite vers
$Y$ donne, étage par étage, des applications compatibles à homotopie près ; la
question est de savoir si l'on peut revenir en arrière, et de façon unique. On
peut toujours trouver une application qui recolle les étages : il suffit de
corriger les représentants un à un, en poussant chaque étage le long d'une
homotopie. L'unicité est une autre affaire : deux applications qui sont
homotopes à chaque étage le sont-elles globalement ? Il faudrait choisir les
homotopies elles-mêmes de façon compatible, et chaque choix peut défaire le
précédent. Les notes s'arrêtent là-dessus, « pas clair du tout », et se
promettent de poser la question à Ronald Brown. La réponse, connue depuis
Milnor (1962), est négative en général : l'obstruction est un groupe
$\varprojlim^{1}$, et il existe des applications non nulles qui sont nulles à
chaque étage.

Les feuillets sont datés par le papier : les listings portent des dates de
juin 1982, ce qui donne une borne inférieure et rien de plus. Le vocabulaire
--- catégorie test, objet de Kan relatif à un localisateur $\underline{W}$ ---
est celui de \emph{Pursuing Stacks}, rédigé à partir de 1983, ce qui s'accorde
avec le titre des archivistes sans le démontrer.

\keywords{presheaves on a product category, external product, weighted limit, cotensor, copower, simplicial objects, simplicial enrichment, matching object, simplicial model category, homotopy extension property, homotopy lifting, test category, sequential homotopy colimit}
\end{resume}

\section*{Le fil du dossier, et les conventions}

\begin{enumerate}
\item[1.] \emph{Préfaisceaux sur un produit} (pages 1 à 3) : les deux lectures
de $(A \times B)^{\wedge}$, les Hom relatifs, les foncteurs d'évaluation
$\Gamma''_{b}$ ; un premier départ en anglais, annulé, pour $B = \Delta$ (page
2).
\item[2.] \emph{Le cotenseur $X^{M}$} (pages 5 et 7) : $\Gamma''_{M}$, limite
pondérée par un préfaisceau $M$ sur $B$, et la stabilité des objets
« représentables en chaque degré ».
\item[3.] \emph{Le tenseur $X \otimes M$ et l'adjonction} (page 9).
\item[4.] \emph{Récapitulation} (page 11).
\item[5.] \emph{Produit extérieur} (pages 13 à 15) : la Proposition et son
Corollaire, à corriger ; l'objet $X^{\dot{\Delta}_{n}}$ (page 14) ;
$p_{B}^{*}$ et les Hom internes (page 15).
\item[6.] \emph{Propriétés universelles de $\widehat{A}$} (page 16, annulée).
\item[7.] \emph{Limites inductives filtrantes dans Hot} (pages 17 à 24) : la
question et la réduction aux monomorphismes (page 18) ; la surjectivité (page
20) ; le carré relatif et sa rectification (pages 21 à 23) ; l'injectivité
(page 24).
\end{enumerate}

Deux suites de lettres de sa main paginent les feuillets : en français,
a à g, sur les pages impaires 3 à 15 ; en anglais, a (page 18), b (page 20),
puis e (page 22) et d (page 24), sans c nulle part. La lecture suit l'ordre des
archivistes, qui est aussi celui de l'argument : le carré relatif des pages 21
à 23 est une digression, et la page 24 s'ouvre sur « Returning to our initial
pb ».\footnote{L'ordre e avant d est celui du dossier ; on ne le corrige pas.
Rien n'indique ce qu'aurait contenu une page c.}

\emph{Conventions.} $A$ et $B$ sont deux petites catégories, $\widehat{A}$ la
catégorie des préfaisceaux d'ensembles sur $A$, dans laquelle on identifie $A$
à son image par Yoneda ; $e_{\widehat{A}}$ est le préfaisceau final et
$\Gamma_{A} = \mathrm{Hom}(e_{\widehat{A}}, -) = \varprojlim_{A^{\circ}}$ le
foncteur des sections. On garde sa notation pour les projections :
\[
p_{B} : A \times B \to A, \qquad p_{A} : A \times B \to B,
\]
l'indice nommant le facteur \emph{oublié}.\footnote{C'est la convention de
toutes les pages, et la seule qui rende ses formules justes : $p_{B*}$ arrive
dans $\widehat{A}$, $p_{A}^{*}(b)$ est l'image réciproque d'un représentable de
$B$. On écrirait aujourd'hui $\mathrm{pr}_{1}$ et $\mathrm{pr}_{2}$.} Pour un
préfaisceau $X$ sur $A \times B$, on note
\[
X' : A^{\circ} \to \widehat{B}, \quad a \mapsto X(a, -), \qquad
X'' : B^{\circ} \to \widehat{A}, \quad b \mapsto X(-, b)
\]
ses deux formes curryfiées. $\underline{\mathrm{Hom}}$ sans indice est le Hom
interne d'une catégorie de préfaisceaux ; les Hom relatifs sont
\[
\underline{\mathrm{Hom}}_{A}(X, Y) = p_{B*}\, \underline{\mathrm{Hom}}(X, Y)
\in \widehat{A}, \qquad
\underline{\mathrm{Hom}}_{B}(X, Y) = p_{A*}\, \underline{\mathrm{Hom}}(X, Y)
\in \widehat{B}.
\]
Pour un préfaisceau $P$, $\pi_{0}(P) = \varinjlim P$ est l'ensemble des
composantes connexes de sa catégorie des éléments, et $P$ est dit
\emph{connexe} si $\pi_{0}(P)$ est un point.\footnote{Cette notation est la
nôtre ; elle sert à énoncer la forme correcte de la Proposition de la page 13
et de ses conséquences.}

\emph{Annoncé, non établi} : la réponse à la question
de la page 18 ; le sens de la flèche $\underline{\mathrm{Hom}}(M, a)
\Rightarrow a$ marquée « ?? » (page 14) ; la proposition de Quillen invoquée
à la page 20, citée et non reproduite.

\pagerange{1}{3}
\section*{Préfaisceaux sur un produit (pages 1 à 3)}

La couverture de la page 1 annonce le programme : « Formalisme des
$X \otimes M$, $X^{M}$, $\underline{\mathrm{Hom}}_{A}$,
$\underline{\mathrm{Hom}}_{B}$ relatif ».\footnote{Le titre est écrit à
l'encre par-dessus l'en-tête d'un listing daté « 06/04/82 » ; il continue par
deux mots lus avec doute, « définitions » et, souligné, « produits ». Le
listing de la page 19, du même travail, est daté « 04 JUIN 82 » : les
feuillets ne sont pas antérieurs à juin 1982.}

\pagerange{2}{2}
\subsection*{Le cas simplicial, premier départ (page 2)}

La page 2, en anglais, annulée par deux traits, part du cas $B = \Delta$. Elle
note $p_{A*}X = (\Gamma_{A} X_{n})_{n}$ et $p_{\Delta*}X = X_{0}$ --- la
seconde formule parce que $[0]$ est final dans $\Delta$, de sorte que $\Gamma_{\Delta}$ est l'évaluation en $[0]$ --- et affirme que, pour un préfaisceau $L$
sur $A$ vu comme objet simplicial constant $p_{\Delta}^{*}(L)$,
\[
\bigl(\underline{\mathrm{Hom}}(L, X_{n})\bigr)_{n} \simeq
\underline{\mathrm{Hom}}\bigl(p_{\Delta}^{*}(L), X\bigr),
\qquad
\bigl(\mathrm{Hom}(L, X_{n})\bigr)_{n} \simeq
\underline{\mathrm{Hom}}_{\Delta}\bigl(p_{\Delta}^{*}(L), X\bigr),
\]
la seconde se déduisant de la première par $p_{A*}$. Les deux sont exactes :
la valeur en $(a, n)$ du second membre de la première est
$\mathrm{Hom}\bigl(p_{\Delta}^{*}(a \times L) \times p_{A}^{*}(\Delta_{n}),
X\bigr) = \mathrm{Hom}(a \times L, X_{n})$. La seconde est l'ensemble
simplicial des applications de $L$ dans $X$, celui qui fait des objets
simpliciaux une catégorie simpliciale.\footnote{La page écrit
$\underline{\mathrm{Hom}}_{A}(L, X_{n})$ pour le Hom interne de $\widehat{A}$,
avec le même indice que pour le Hom relatif $\underline{\mathrm{Hom}}_{A} =
p_{\Delta*}\,\underline{\mathrm{Hom}}$ qu'elle vient de définir ; on supprime
l'indice. La vérification est la nôtre : la phrase qui devait la donner
s'arrête sur « This is easily checked (see below) to be », « justified » étant
ajouté au-dessus. Le sommet de gauche du premier diagramme est lu avec doute ;
la marge gauche porte, en travers, des fragments sans lien apparent avec le
texte ($X(a, b) = M(b)$, $u \mapsto u(i_{0})$, …).}

\pagerange{3}{3}
\subsection*{Deux lectures d'un même préfaisceau (page 3)}

On a les identifications
\[
(A \times B)^{\wedge} \simeq \underline{\mathrm{Hom}}(A^{\circ}, \widehat{B})
\simeq \underline{\mathrm{Hom}}(B^{\circ}, \widehat{A}),
\qquad X \mapsto X', \ X \mapsto X'',
\]
et les foncteurs $p_{B*} : (A \times B)^{\wedge} \to \widehat{A}$,
$p_{A*} : (A \times B)^{\wedge} \to \widehat{B}$, adjoints à droite de
$p_{B}^{*}$ et $p_{A}^{*}$. Comme $q_{A} p_{B} = q_{B} p_{A}$ est la flèche vers
la catégorie ponctuelle,
\[
\mathrm{Hom}(X, Y) \simeq \Gamma_{A}\, \underline{\mathrm{Hom}}_{A}(X, Y)
\simeq \Gamma_{B}\, \underline{\mathrm{Hom}}_{B}(X, Y).
\]

\emph{Évaluation en $b$.} Pour $b \in \mathrm{Ob}\, B$, soit
$\Gamma''_{b} : (A \times B)^{\wedge} \to \widehat{A}$, $X \mapsto X''(b) =
X(-, b)$. Alors, fonctoriellement en $b$ et $X$,
\[
\Gamma''_{b}(X) \simeq p_{B*}\, \underline{\mathrm{Hom}}\bigl(p_{A}^{*}(b),
X\bigr) = \underline{\mathrm{Hom}}_{A}\bigl(p_{A}^{*}(b), X\bigr).
\]
En effet, pour $a \in \mathrm{Ob}\, A$, l'adjonction $p_{B}^{*} \dashv p_{B*}$
puis celle du Hom interne donnent
\[
\mathrm{Hom}_{\widehat{A}}\bigl(a, p_{B*}\, \underline{\mathrm{Hom}}(p_{A}^{*}(b),
X)\bigr) \simeq \mathrm{Hom}\bigl(p_{B}^{*}(a) \times p_{A}^{*}(b), X\bigr)
= X(a, b),
\]
parce que $p_{B}^{*}(a) \times p_{A}^{*}(b)$ est le préfaisceau représenté par
$(a, b)$.\footnote{La page termine la chaîne par « OK ». Plusieurs indices
$A$, $B$ y sont récrits l'un sur l'autre ; on suit la lecture finale, qui est
la seule cohérente avec la convention des projections. Au-dessus, la page dit
$\Gamma''_{b}$ « déduit par $(A \times B)^{\wedge} \simeq
\underline{\mathrm{Hom}}(A^{\circ}, B)$ » : il faut lire $\widehat{B}$. Un
paragraphe symétrique, sur les $\Gamma'_{a} = X'(a) \simeq p_{A*}\,
\underline{\mathrm{Hom}}(p_{B}^{*}(a), X)$, est annulé par trois traits ; il
est juste. Les mots de la phrase qui introduit $\Gamma''_{b}$ sont en partie
illisibles ; on ne les restitue pas.}

\pagerange{5}{7}
\section*{Le cotenseur $X^{M}$ (pages 5 et 7)}

\emph{Une troisième lecture de $\Gamma''_{b}$.} Soit $\varepsilon_{b} :
\widehat{B} \to (\mathrm{Ens})$, $G \mapsto \mathrm{Hom}(b, G) = G(b)$ ; alors
$\Gamma''_{b}(X) = \varepsilon_{b} \circ X'$. Cela se généralise à un
préfaisceau quelconque $M$ sur $B$ : avec $\varepsilon_{M}(G) =
\mathrm{Hom}_{\widehat{B}}(M, G)$, posons
\[
\Gamma''_{M}(X) = \varepsilon_{M} \circ X' : \quad a \longmapsto
\mathrm{Hom}_{\widehat{B}}\bigl(M, X(a, -)\bigr).
\]
C'est la limite du diagramme $X'' : B^{\circ} \to \widehat{A}$ pondérée par
$M$.\footnote{Le nom de limite pondérée est le nôtre ; il n'est pas sur la
page. La page écrit ici $X'' : A^{\circ} \to \widehat{B}$, alors que la page
3 réserve $X''$ à la forme $B^{\circ} \to \widehat{A}$, et
« $X'' \mapsto \varepsilon_{M} \circ \Gamma''_{M}$ » pour la définition, où le
sens demande $\varepsilon_{M} \circ X'$. La première ligne de la page est
biffée et récrite au-dessus ; la nouvelle lecture est douteuse.} On a,
fonctoriellement en $M$ et $X$,
\[
\Gamma''_{M}(X) \simeq p_{B*}\, \underline{\mathrm{Hom}}\bigl(p_{A}^{*}(M),
X\bigr),
\]
car les deux membres, évalués en $a$, valent
$\mathrm{Hom}\bigl(p_{B}^{*}(a) \times p_{A}^{*}(M), X\bigr)$ : à droite par
les deux adjonctions comme à la page 3, à gauche parce que
$\mathrm{Hom}_{\widehat{B}}(M, X'(a)) = \mathrm{Hom}_{\widehat{B}}(M,
p_{A*}\underline{\mathrm{Hom}}(p_{B}^{*}(a), X))$.\footnote{La page conduit les
deux calculs en colonnes et les relie par une double flèche courbe marquée
« OK ! ». Le terme $p_{A*}(\underline{\mathrm{Hom}}(p_{B}^{*}(a), X))$ de la
colonne de gauche est raturé et de lecture douteuse ; c'est $\Gamma'_{a}(X) =
X'(a)$, et c'est ce qu'il doit être.}

\emph{Le cotenseur.} On pose, pour $M \in \mathrm{Ob}\, \widehat{B}$,
\[
X^{M} \overset{\text{déf}}{=} \underline{\mathrm{Hom}}\bigl(p_{A}^{*}(M),
X\bigr) \in (A \times B)^{\wedge},
\qquad \text{d'où} \qquad \Gamma''_{M}(X) = p_{B*}\, X^{M}.
\]
Sa valeur en $(a, b)$ est $\mathrm{Hom}_{\widehat{B}}\bigl(b \times M, X(a,
-)\bigr)$ ; autrement dit $(X^{M})'' (b) = \Gamma''_{b \times M}(X)$.

\emph{Stabilité} (page 7). On s'intéresse à la sous-catégorie pleine
$\underline{\mathrm{Hom}}(B^{\circ}, A)$ de $(A \times B)^{\wedge}$, formée des
$X$ dont chaque $X''(b)$ est représentable --- pour $B = \Delta$, la catégorie
$sA$ des objets simpliciaux de $A$. Si $B$ a un objet final $e_{B}$, $p_{B*}$
est l'évaluation en $e_{B}$, $p_{B*} = \Gamma''_{e_{B}}$, et envoie donc
$\underline{\mathrm{Hom}}(B^{\circ}, A)$ dans $A$.\footnote{La page écrit
« $p_{B*}$ $(= (\Gamma_{B})^{\circ})$ » suivi d'un mot illisible, qu'on lit
par le sens : $p_{B*}$ envoie $\underline{\mathrm{Hom}}(B^{\circ}, A)$ dans $A$.
L'objet final est supposé à cet endroit, et la justification « car il
s'identifie alors à $\Gamma''_{e_{B}}$ » est de sa main. Le diagramme qui
suit a deux flèches qui convergent vers un point où rien n'est écrit.}
$X^{M}$ reste dans $\underline{\mathrm{Hom}}(B^{\circ}, A)$ dès que $A$ a les
limites projectives requises : $(X^{M})''(b)$ est la limite de $X''$ pondérée
par $b \times M$, c'est-à-dire la limite dans $A$ du diagramme des
$X''(c)$ indexé par la catégorie des éléments de $b \times M$. Pour $B =
\Delta$ et $M$ fini, des limites finies suffisent.\footnote{Il écrit :
« moyennant éventuellement des conditions de finitude », et « il suffit que $A$
soit stable par limites du type requis pour représenter $M$ » ; la précision
sur la catégorie des éléments de $b \times M$ est la nôtre. Le mot
« limites » est lu avec doute.} En revanche,
$\underline{\mathrm{Hom}}(B^{\circ}, A)$ ne contient pas les $p_{A}^{*}(M)$
eux-mêmes, sauf cas dégénéré : il faudrait que, pour tout $b$, le préfaisceau
constant de valeur $M(b)$ sur $A$ soit représentable, ce qui force $M(b)$ à
être un point --- un représentable $h_{x}$ constant a $\mathrm{End}(x)$ réduit
à l'identité --- et $A$ à avoir un objet final : $M = e_{\widehat{B}}$.
« !! »\footnote{C'est une note de bas de page de sa main, serrée sur trois
lignes, dont plusieurs mots sont illisibles ; sa conclusion, « $M$ objet
final de $B^{\wedge}$, et $A$ admet objet final !! », se lit. L'argument entre
tirets est le nôtre.}

Enfin, pour $X, Y$ dans $(A \times B)^{\wedge}$, en appliquant $p_{A*}$ à
l'adjonction du Hom interne,
\[
\underline{\mathrm{Hom}}_{B}(X, Y^{M}) \simeq
\underline{\mathrm{Hom}}_{B}\bigl(X \times p_{A}^{*}(M), Y\bigr).
\]

\pagerange{9}{9}
\section*{Le tenseur $X \otimes M$ (page 9)}

Pour $X, Y$ dans $\underline{\mathrm{Hom}}(B^{\circ}, A)$, on voudrait
trouver le membre de gauche de cette formule dans
$\underline{\mathrm{Hom}}(B^{\circ}, A)$ aussi. Or $X \times p_{A}^{*}(M)$ est,
comme foncteur $B^{\circ} \to \widehat{A}$,
\[
b \longmapsto X''(b) \times M(b) = \coprod^{\widehat{A}}_{\sigma \in M(b)}
X''(b),
\]
une somme dans $\widehat{A}$ de représentables, qui n'est presque jamais
représentable. Supposons que $A$ ait des sommes (des sommes finies si les
$M(b)$ sont finis), et définissons $X \otimes M \in
\underline{\mathrm{Hom}}(B^{\circ}, A)$ par
\[
(X \otimes M)(b) = \coprod^{A}_{\sigma \in M(b)} X''(b),
\]
la somme étant prise cette fois dans $A$. La flèche canonique va de la somme
dans $\widehat{A}$ vers la somme dans $A$,
\[
i_{b} : \bigl(X \times p_{A}^{*}(M)\bigr)''(b) =
\coprod^{\widehat{A}}_{\sigma \in M(b)} X''(b) \longrightarrow
\coprod^{A}_{\sigma \in M(b)} X''(b),
\]
et elle induit une bijection sur les flèches vers un représentable ; c'est la
définition même de la somme dans $A$.\footnote{La page trace la flèche
$i_{b}$ (lecture douteuse de l'indice) de la somme dans $A$ vers la somme dans
$\widehat{A}$ ; il n'y a pas de flèche naturelle dans ce sens. Les deux sommes
portent en exposant la catégorie où elles sont prises, comme ici.} D'où, pour
$X, Y$ dans $\underline{\mathrm{Hom}}(B^{\circ}, A)$, fonctoriellement,
\[
\underline{\mathrm{Hom}}_{B}(X, Y^{M}) \simeq \underline{\mathrm{Hom}}_{B}(X
\otimes M, Y).
\]
Le calcul est le suivant. La valeur en $b'$ du membre de droite de la formule
de la page 7 est $\mathrm{Hom}\bigl(X \times p_{A}^{*}(M \times b'),
Y\bigr)$ ; comme les $Y''(c)$ sont représentables, une flèche
$X \times p_{A}^{*}(N) \to Y$ est la même chose qu'une flèche $X \otimes N \to
Y$, et $(X \otimes M) \otimes b' = X \otimes (M \times b')$.\footnote{Cette
vérification est la nôtre ; la page énonce la formule. Sous l'accolade de la
formule précédente, une notation raturée laisse deviner $X \otimes M$.}

Pour $B = \Delta$, c'est la structure simpliciale de la catégorie $sA$ des
objets simpliciaux de $A$ : $(X \otimes K)_{n} = \coprod_{K_{n}} X_{n}$ pour
un ensemble simplicial $K$, l'ensemble simplicial
$\underline{\mathrm{Hom}}_{\Delta}(X, Y)$ des applications de $X$ dans $Y$, et
le cotenseur $Y^{K}$ ; $sA$ est ainsi \emph{tensorisée} et
\emph{cotensorisée} sur les ensembles simpliciaux.\footnote{Les deux mots sont
les nôtres. C'est la structure qu'utilise Quillen (\emph{Homotopical Algebra},
1967, chap.~II, §1--2) pour les objets simpliciaux d'une catégorie ayant les
limites et les sommes finies, et c'est elle que la seconde partie du dossier
emploie sous la forme du cylindre $X \otimes J$.}

\pagerange{11}{11}
\section*{Récapitulation (page 11)}

Notons ${}^{g}L = p_{B}^{*}(L)$ pour $L \in \mathrm{Ob}\, \widehat{A}$ et
${}^{d}M = p_{A}^{*}(M)$ pour $M \in \mathrm{Ob}\, \widehat{B}$, et $L
\boxtimes M = {}^{g}L \times {}^{d}M$ ; pour $a$, $b$ représentables, $a
\boxtimes b$ est le représentable $(a, b)$. Alors
\[
X(L, M) \overset{\text{déf}}{=} \mathrm{Hom}(L \boxtimes M, X), \qquad X(a, b)
= \Gamma_{A \times B}\, \underline{\mathrm{Hom}}(a \boxtimes b, X),
\]
\[
{}^{g}X(L) = \underline{\mathrm{Hom}}_{B}({}^{g}L, X) \in \widehat{B}, \qquad
{}^{d}X(M) = \underline{\mathrm{Hom}}_{A}({}^{d}M, X) = p_{B*}\, X^{M} =
\Gamma''_{M}(X) \in \widehat{A},
\]
avec, pour $L = a$, ${}^{g}X(a) = X'(a)$, et pour $M = b$, ${}^{d}X(b) =
X''(b)$ ; et l'adjonction de la page 9, $\underline{\mathrm{Hom}}_{B}(X,
Y^{M}) \simeq \underline{\mathrm{Hom}}_{B}(X \otimes M, Y)$.\footnote{La
première ligne de la page écrit $p_{B*}$ dans les formules de ${}^{g}X(a)$ et
de ${}^{d}X(b)$, l'indice du premier étant surchargé ; pour ${}^{g}X(a)$ il
faut $p_{A*}$, comme la page l'écrit d'ailleurs deux lignes plus bas pour
${}^{g}X(L)$. L'accolade marquée $X^{M}$ couvre
$p_{B*}\,\underline{\mathrm{Hom}}({}^{d}M, X)$, qui est $p_{B*}X^{M}$ et non
$X^{M}$ lui-même. Les mots « identifié à » et « identifiés », au-dessus de
$p_{B}^{*}(a)$ et $p_{A}^{*}(b)$, sont lus avec doute.}

\pagerange{13}{13}
\section*{Le produit extérieur (page 13)}

Les identifications de la page 3 se prolongent aux préfaisceaux :
\[
(A \times B)^{\wedge} \simeq \underline{\mathrm{Hom}}(A^{\circ}, \widehat{B})
\simeq \underline{\mathrm{Hom}}^{!}(\widehat{A}^{\circ}, \widehat{B})
\simeq \underline{\mathrm{Hom}}^{!!}\bigl(\widehat{A}^{\circ} \times
\widehat{B}^{\circ}, (\mathrm{Ens})\bigr),
\]
où $\underline{\mathrm{Hom}}^{!}$ désigne les foncteurs qui commutent aux
limites projectives quelconques --- ceux qui transforment les limites
inductives de $\widehat{A}$ en limites projectives --- et
$\underline{\mathrm{Hom}}^{!!}$ ceux qui le font en chaque variable
séparément. Le préfaisceau $X$ correspond au foncteur $(L, M) \mapsto X(L, M)
= \mathrm{Hom}(L \boxtimes M, X)$.\footnote{La page note ces conditions par
un « ! » et un « !! » en indice, et la marge précise : « le signe ! signifie :
foncteurs commutant aux $\varprojlim$ quelconques ! » ; la page 16 met le
signe en indice pour la commutation aux limites inductives et en exposant
pour les limites projectives, et c'est cette dernière convention qu'on suit
partout. La ligne se termine par un dernier terme,
$\underline{\mathrm{Hom}}_{!}((\widehat{A} \times \widehat{B})^{\circ},
\mathrm{Ens})$, avec un seul signe ; pris au sens d'une condition sur le
produit, il décrirait les préfaisceaux sur la somme $A \sqcup B$,
c'est-à-dire $\widehat{A} \times \widehat{B}$, et non $(A \times
B)^{\wedge}$. C'est la condition séparée !! qui convient ; la fin de la note
marginale, « et un seul signe ! relatif », est en partie illisible.}

\textbf{Proposition} (forme exacte). Pour $L, L' \in \mathrm{Ob}\,
\widehat{A}$ et $M, M' \in \mathrm{Ob}\, \widehat{B}$,
\[
\mathrm{Hom}(L' \boxtimes M', L \boxtimes M) \simeq
\mathrm{Hom}(L', L)^{\pi_{0}(M')} \times \mathrm{Hom}(M', M)^{\pi_{0}(L')}.
\]
En particulier, si $L'$ et $M'$ sont connexes,
\[
\mathrm{Hom}(L' \boxtimes M', L \boxtimes M) \simeq \mathrm{Hom}(L', L)
\times \mathrm{Hom}(M', M),
\]
et le foncteur $(L, M) \mapsto L \boxtimes M$, de $\widehat{A} \times
\widehat{B}$ dans $(A \times B)^{\wedge}$, est pleinement fidèle sur les
couples de préfaisceaux connexes.

\emph{Démonstration.} Un morphisme vers un produit est un couple de
morphismes, et $\mathrm{Hom}(L' \boxtimes M', p_{B}^{*}(L)) \simeq
\mathrm{Hom}(p_{B!}(L' \boxtimes M'), L)$ avec $p_{B!}(L' \boxtimes M') = L'
\times \pi_{0}(M')$, la somme de $\pi_{0}(M')$ copies de $L'$ ; de même pour
l'autre facteur.

La page énonce la Proposition sans restriction : $(L, M) \mapsto L \boxtimes
M$ pleinement fidèle, et $\mathrm{Hom}(L' \boxtimes M', L \boxtimes M) =
\mathrm{Hom}(L', L) \times \mathrm{Hom}(M', M)$. C'est faux en général : pour
$M' = \varnothing$, le premier membre est un point et le second est
$\mathrm{Hom}(L', L)$.\footnote{La démonstration de la page vérifie la formule
pour $L' = a$, $M' = b$ représentables, où elle est vraie « par définition »,
puis observe que les deux membres, comme foncteurs de $(L', M')$,
transforment limites inductives en limites projectives dans les deux
variables. C'est vrai du premier membre, puisque $\boxtimes$ commute aux
limites inductives en chaque variable, et faux du second : $\mathrm{Hom}(L',
L) \times \mathrm{Hom}(M', M)$, à $M'$ fixé, ne transforme en limites
projectives que les limites inductives \emph{connexes} en $L'$, le facteur
constant faisant obstacle aux sommes et à la limite vide. Un préfaisceau
connexe étant limite inductive connexe de représentables (sa catégorie des
éléments est connexe), l'argument établit exactement la forme restreinte. Il
écrit d'abord « Corollaire », biffé, puis « Proposition ». Il fait aussi
partir le foncteur de $(\widehat{A} \times \widehat{B})^{\circ}$ ; la variance
correcte est covariante, de $\widehat{A} \times \widehat{B}$, et c'est ce
qu'il utilise en identifiant $L \boxtimes M$ au foncteur qu'il représente
sur $\widehat{A} \times \widehat{B}$. La forme générale avec les exposants
$\pi_{0}$ est la nôtre.}

\textbf{Corollaire.} Supposons que, pour tous $a \in \mathrm{Ob}\, A$ et $b
\in \mathrm{Ob}\, B$, les préfaisceaux $a \times L'$ et $b \times M'$ soient
connexes. Alors
\[
\underline{\mathrm{Hom}}(L' \boxtimes M', L \boxtimes M) \simeq
\underline{\mathrm{Hom}}(L', L) \boxtimes \underline{\mathrm{Hom}}(M', M).
\]
On évalue les deux membres en $(a, b)$, en utilisant $(a \boxtimes b) \times
(L' \boxtimes M') = (a \times L') \boxtimes (b \times M')$ et la
Proposition.\footnote{La page démontre le Corollaire ainsi, sans hypothèse,
et note « ok » au pied du calcul ; l'hypothèse de connexité est celle que la
Proposition exige pour $a \times L'$ et $b \times M'$, et sans elle la formule
peut être fausse (prendre $A$ discret à deux objets et $L'$ représenté par l'un
d'eux). Elle est satisfaite si $L'$ et $M'$ sont connexes et si les produits de
deux représentables de $\widehat{A}$, et de $\widehat{B}$, sont connexes ---
par exemple si $A$ et $B$ ont des produits, ou pour $B = \Delta$, où
$\Delta_{n} \times \Delta_{m}$ est connexe. Sur la page, les deux colonnes du
calcul sont côte à côte.}

\pagerange{14}{14}
\section*{Les objets $X^{\Delta_{n}}$ et $X^{\dot{\Delta}_{n}}$ (page 14)}

Une page de formules isolées, pour $B = \Delta$ et $X$ un objet simplicial :
\[
(X^{\Delta_{n}})_{0} = p_{\Delta*}\bigl(\underline{\mathrm{Hom}}(p_{A}^{*}(\Delta_{n}),
X)\bigr) = X_{n}, \qquad
(X^{\dot{\Delta}_{n}})_{0} = \varprojlim_{\Delta_{k} \to \dot{\Delta}_{n}} X_{k}
= (\mathrm{Cosk}_{n-1} X)_{n},
\]
la seconde étant la limite des $X_{k}$ sur les faces propres du simplexe ---
ce qu'on appelle aujourd'hui l'\emph{objet de recollement} (matching object)
de $X$ en degré $n$.\footnote{Le nom est le nôtre. La page écrit la limite
« $\varprojlim(X_{n-1}, X_{n-2}, \cdots)$ », arguments lus avec doute, et
porte en marge $(\mathrm{Cosk}_{n-1}X)_{n}$, ainsi qu'un petit triangle
hachuré dont les côtés sont marqués $x_{1}$ ; $\dot{\Delta}_{n}$ est le bord
de $\Delta_{n}$.} Puis, pour un morphisme $X \to Y$, les deux flèches
\[
X^{\Delta_{n}} \longrightarrow X^{\dot{\Delta}_{n}}
\times_{Y^{\dot{\Delta}_{n}}} Y^{\Delta_{n}}, \qquad
X_{n} \longrightarrow (X^{\dot{\Delta}_{n}})_{0}
\times_{(Y^{\dot{\Delta}_{n}})_{0}} Y_{n},
\]
la seconde étant la composante de degré $0$ de la première. Ce sont les
flèches dont Quillen demande qu'elles soient des fibrations dans l'axiome
des catégories de modèles simpliciales, et dont la seconde définit
aujourd'hui les fibrations de Reedy.\footnote{Les deux rapprochements sont les
nôtres ; la page ne porte aucun texte.}

Restent trois formules sans texte, dont une question : pour $a$ un objet de
$A$ vu comme objet simplicial constant, a-t-on $\underline{\mathrm{Hom}}(M, a)
\simeq a$, c'est-à-dire $a^{M} = a$ ? La page la marque « ?? » ; la page 15
y répond, sous une hypothèse sur $M$ qu'elle ne dit pas.\footnote{Les deux
autres formules, $\mathrm{Hom}(X \times M, a) \simeq \mathrm{Hom}(X, -)$ et
$\mathrm{Hom}(\alpha \boxtimes (\beta \times M), \alpha \boxtimes e_{B}) \simeq
\mathrm{Hom}(\alpha \boxtimes \beta, \alpha \boxtimes e_{B})$, sont
inachevées, la seconde touchant le bord de la feuille et de lecture douteuse
dans son dernier argument. Elles préparent le calcul de la page 15.}

\pagerange{15}{15}
\section*{$p_{B}^{*}$ et les Hom internes (page 15)}

\textbf{Corollaire.} $p_{B}^{*} : \widehat{A} \to (A \times B)^{\wedge}$
commute aux Hom internes,
\[
\underline{\mathrm{Hom}}\bigl(p_{B}^{*}(L'), p_{B}^{*}(L)\bigr) \simeq
p_{B}^{*}\, \underline{\mathrm{Hom}}(L', L),
\]
et de même $p_{A}^{*}$. C'est le Corollaire de la page 13 pour $M' = M =
e_{\widehat{B}}$, où son hypothèse n'est pas nécessaire : avec $M =
e_{\widehat{B}}$, le second facteur de la formule exacte est un point et le
premier vaut $\mathrm{Hom}(a \times L', L)^{\pi_{0}(b)} = \mathrm{Hom}(a
\times L', L)$. C'est aussi un cas de la formule de projection :
$p_{B!}(F \times p_{B}^{*}L) = p_{B!}(F) \times L$, parce que le produit par un
ensemble fixe commute aux limites inductives d'ensembles.\footnote{La page
écrit « Il suffit de prendre $M' = M = e_{B^{\wedge}}$ », la première
tentative, $M' \simeq e_{B^{\wedge}}$, étant biffée ; puis, dans la formule,
garde $M'$ quelconque :
$\underline{\mathrm{Hom}}(L' \boxtimes M', L \boxtimes e_{\widehat{B}}) \simeq
\underline{\mathrm{Hom}}(L', L) \boxtimes e_{\widehat{B}}$. Sous cette forme
la formule est fausse ($M' = \varnothing$) ; elle demande que $b \times M'$
soit connexe pour tout $b$. La justification par la formule de projection est
la nôtre.}

\emph{Le cotenseur d'un objet constant.} Identifions $M \in \mathrm{Ob}\,
\widehat{B}$ à $e_{\widehat{A}} \boxtimes M = p_{A}^{*}(M)$ et $L \in \mathrm{Ob}\,
\widehat{A}$ à $L \boxtimes e_{\widehat{B}} = p_{B}^{*}(L)$. Il existe
toujours une flèche canonique
\[
L \longrightarrow L^{M},
\]
cas de la flèche $Q \to \underline{\mathrm{Hom}}(P, Q)$ de toute catégorie
cartésienne fermée, et
\[
L^{M}(a, b) = L(a)^{\pi_{0}(b \times M)} .
\]
Elle est donc un isomorphisme si et seulement si $b \times M$ est connexe pour
tout $b$ ; pour $B = \Delta$, cela revient à dire que $M$ est un ensemble
simplicial connexe, puisque $\pi_{0}$ commute aux produits finis d'ensembles
simpliciaux. C'est la réponse à la question « ?? » de la page 14.\footnote{La
page conclut « $L^{M} = L$ pour $L \in A^{\wedge}$, $M \in \mathrm{Ob}\,
B^{\wedge}$ », sans hypothèse, en faisant $L' = e_{A^{\wedge}}$ et en changeant
$M'$ en $M$ dans la formule précédente, que l'on vient de voir fausse pour
$M'$ quelconque. Pour $M = \varnothing$, $L^{\varnothing}$ est le préfaisceau
final et non $L$. L'expression de $L^{M}(a, b)$ et le critère sont les nôtres.
Les mots qui précèdent « flèche canonique », et deux lignes biffées dessous,
la première lue « produits », sont de lecture douteuse.}

\pagerange{16}{16}
\section*{Les propriétés universelles de $\widehat{A}$ (page 16)}

Une page annulée par un long trait oblique, où $M$ désigne cette fois une
catégorie, rappelle la propriété universelle de $\widehat{A}$ et sa duale.
Notons $\underline{\mathrm{Hom}}_{!}$ les foncteurs qui commutent aux limites
inductives, $\underline{\mathrm{Hom}}^{!}$ ceux qui commutent aux limites
projectives, et $A^{\vee} = (A^{\circ})^{\wedge}$. Si $M$ a des limites
inductives,
\[
\underline{\mathrm{Hom}}(A, M) \simeq \underline{\mathrm{Hom}}_{!}(\widehat{A},
M), \qquad
\underline{\mathrm{Hom}}(A^{\circ}, M) \simeq
\underline{\mathrm{Hom}}_{!}(A^{\vee}, M) ;
\]
si $M$ a des limites projectives, $M^{\circ}$ a des limites inductives, et
\[
\underline{\mathrm{Hom}}(A, M^{\circ}) \simeq
\underline{\mathrm{Hom}}_{!}(\widehat{A}, M^{\circ}), \qquad\text{soit}\qquad
\underline{\mathrm{Hom}}(A^{\circ}, M) \simeq
\underline{\mathrm{Hom}}^{!}(\widehat{A}^{\circ}, M).
\]
$\widehat{A}$ est la complétion de $A$ par limites inductives, et
$\widehat{A}^{\circ}$ celle de $A^{\circ}$ par limites projectives. Le second
cas, pour $M = \widehat{B}$, est l'identification $(A \times B)^{\wedge}
\simeq \underline{\mathrm{Hom}}^{!}(\widehat{A}^{\circ}, \widehat{B})$ de la
page 13.\footnote{Les noms de complétion sont les nôtres. La page écrit sous
$M$ « stable par $\varinjlim$ » et « catégorie stable par $\varprojlim$ »,
lectures douteuses. Une formule encadrée est barrée, et une autre,
$\underline{\mathrm{Hom}}_{!}(A^{\vee}, M)^{\circ} \simeq
\underline{\mathrm{Hom}}_{!}(A^{1}, M^{\circ})$, porte un exposant lu avec
doute « 1 » qu'on n'interprète pas.}

\pagerange{17}{18}
\section*{Limites inductives filtrantes dans Hot (pages 17 et 18)}

La couverture de la page 17, « Question des $\varinjlim$ filtrantes dans
Hot », ouvre la seconde partie, en anglais.\footnote{Titre seul sur une
feuille vierge. Les pages 18 à 24 sont en anglais ; on les lit en
français.}

\emph{La question.} La catégorie homotopique $\mathrm{Hot}$ a-t-elle des
limites inductives filtrantes, et les foncteurs canoniques
\[
(\mathrm{Cat}) \longrightarrow \mathrm{Hot}(\underline{W}), \qquad
\widehat{A} \longrightarrow \mathrm{Hot}_{A} \simeq \mathrm{Hot}(\underline{W})
\quad (A \text{ catégorie test})
\]
y commutent-ils ? Le premier cas se ramène au second pour $A = \Delta$ : le
foncteur nerf $(\mathrm{Cat}) \to \widehat{\Delta}$, défini par le foncteur
test standard $i : \Delta \to (\mathrm{Cat})$, commute aux limites inductives
filtrantes, parce que les catégories $i([n])$ sont de présentation finie.\footnote{Le
vocabulaire --- catégorie test, $\mathrm{Hot}_{A}$, localisateur
$\underline{W}$ --- est celui de \emph{Pursuing Stacks}. La parenthèse sur
« of f.p. » n'est pas refermée sur la page.}

\emph{Réduction aux monomorphismes.} Pour un système indexé par un ensemble
dénombrable, on se ramène à une suite $X_{0} \xrightarrow{f_{0}} X_{1}
\xrightarrow{f_{1}} \cdots$. En factorisant chaque flèche en un monomorphisme
suivi d'une équivalence d'homotopie, on construit par récurrence une suite
$X'_{0} \hookrightarrow X'_{1} \hookrightarrow \cdots$ de monomorphismes et des
équivalences $u_{i} : X'_{i} \to X_{i}$ compatibles, $u_{0} = \mathrm{id}$,
d'où $u_{\infty} : X'_{\infty} \to X_{\infty}$ entre les limites inductives.
Il faut alors voir (i) que $X'_{\infty}$ est limite inductive des $X'_{i}$
dans $\mathrm{Hot}_{A}$, et (ii) que $u_{\infty}$ est dans
$\underline{W}$. Le point (ii) est vrai : les équivalences faibles sont
stables par limites inductives filtrantes.\footnote{La page ne tranche pas
(ii) ; elle le dit équivalent, une fois (i) acquis, à la commutation
cherchée, l'expression « is equivalent to » étant biffée puis prolongée d'une
flèche, et la fin de la phrase de lecture douteuse. La stabilité est
classique pour les ensembles simpliciaux ; pour une catégorie test $A$, elle
s'en déduit parce que le foncteur « catégorie des éléments » $\widehat{A}
\to (\mathrm{Cat})$ est un adjoint à gauche et que le nerf commute aux limites
inductives filtrantes. L'argument est le nôtre.} Tout repose donc sur (i),
c'est-à-dire sur la propriété
\[
(*) \qquad \mathrm{Hom}_{\mathrm{Hot}}(X_{\infty}, Y) \longrightarrow
\varprojlim_{i} \mathrm{Hom}_{\mathrm{Hot}}(X_{i}, Y)
\quad \text{bijective,}
\]
pour une suite de monomorphismes. On peut supposer $Y$ de Kan relativement à
$\underline{W}$, c'est-à-dire fibrant ; alors $\mathrm{Hom}_{\mathrm{Hot}}(X,
Y)$ est l'ensemble des classes d'homotopie de flèches $X \to Y$, et la
relation d'homotopie entre flèches $X \to Y$ coïncide avec l'homotopie
élémentaire.\footnote{Ces deux affirmations ouvrent la page 20. Elles sont
justes dans la structure de modèles de Kan--Quillen sur les ensembles
simpliciaux, où tout objet est cofibrant ; « the », « among » et « the same »
sont ajoutés entre les lignes.}

\pagerange{20}{20}
\section*{La surjectivité de $(*)$ (page 20)}

Soit une famille $g_{n} : X_{n} \to Y$ avec $g_{n+1} f_{n} \sim g_{n}$ pour
tout $n$. Il s'agit de trouver $g'_{n} \sim g_{n}$ avec $g'_{n+1} f_{n} =
g'_{n}$ : les $g'_{n}$ définissent alors une flèche $X_{\infty} \to Y$ qui
relève la famille. C'est la \emph{surjectivité} de $(*)$.\footnote{La page
introduit ce problème par « Injectivity of $(*)$? ». Le problème posé
--- remplacer une famille compatible à homotopie près par une famille
strictement compatible --- est celui de la surjectivité, et la page 24 le
dit : « surjection (using strongly that the transition morphisms are
cofibrations) », avant de passer à l'injectivité.} Il se place dans une
situation plus générale : un carré

\begin{tikzcd}[column sep=large, row sep=large]
X \arrow[r, "g"] \arrow[d, "i"'] & Y \arrow[d, "p"] \\
X' \arrow[r, "g'"'] \arrow[ur, dashed, "\varphi"] & Y'
\end{tikzcd}

où $p g = g' i$, $i$ est une cofibration, $p$ une fibration, $X$, $X'$
cofibrants et $Y$, $Y'$ fibrants, dans une catégorie de modèles quelconque,
et $\varphi$ une flèche rendant les deux triangles commutatifs à homotopie
près, $\varphi i \overset{k}{\sim} g$ et $p \varphi \overset{h}{\sim} g'$.
Existe-t-il $\overline{\varphi} \sim \varphi$ qui les rende commutatifs,
$\overline{\varphi} i = g$ et $p \overline{\varphi} = g'$ ?\footnote{Sous ces
hypothèses, homotopie à gauche et à droite coïncident ; la page le dit, les
deux signes $\sim$ portant des lettres lues avec doute $\ell$ et $r$. Dans la
dernière formule elle écrit $\varphi$ pour $\overline{\varphi}$. Le cas de la
suite est $i = f_{n}$, $p : Y \to e$, $g = g_{n}$, $\varphi = g_{n+1}$, la
restitution de $f_{n} : X_{n} \to X_{n+1}$ étant douteuse. Une phrase biffée
commence la construction d'un objet cylindre.} Il y voit « à peu près
exactement le contenu » de la version de Quillen des théorèmes d'extension
des homotopies pour les catégories de modèles simpliciales (proposition 4 du
chapitre II, n°~2.4), « bien qu'il ait été un peu difficile
d'y trouver un sens, telle qu'il l'énonçait ».\footnote{La référence est celle
de la page ; on ne l'a pas confrontée ici au texte de Quillen.}

Quand $Y' = e$, la réponse est oui, et c'est tout ce que la suite demande : la
flèche $X' \times \{0\} \cup X \times \Delta(1) \to Y$ définie par $\varphi$ et
$k$ s'étend à $X' \times \Delta(1)$, puisque l'inclusion est une cofibration
triviale et $Y$ fibrant ; l'extrémité de l'homotopie étendue est
$\overline{\varphi}$. Par récurrence, $g'_{0} = g_{0}$ et $g'_{n+1}$ obtenue
ainsi à partir de $\varphi = g_{n+1}$ et de $g = g'_{n}$, et $(*)$ est
surjective.\footnote{La récurrence est la nôtre ; la page 22 dit « in case
$Y'$ is the final object, we are through! », et la page 24 tient la
surjectivité pour acquise.}

Dans le cas général, la note marginale dit : c'est possible si et seulement
si l'on peut choisir $k$ et $h$ de sorte que les deux homotopies $p \circ k$
et $h \circ (i \otimes J)$, qui vont toutes deux de $p \varphi i$ à $g' i$,
soient homotopes --- relativement à leurs extrémités. « Ce n'est sans doute
pas toujours satisfait ! » Le critère est exact, et la condition peut
manquer.\footnote{La note est écrite en travers dans la marge gauche, et ses
lectures sont très douteuses ; la précision « relativement à leurs
extrémités » est la nôtre. \emph{Nécessité} : si $\overline{\varphi}$ existe,
une homotopie $K : \varphi \sim \overline{\varphi}$ fournit $k = K \circ (i
\otimes J)$ et $h = p \circ K$, et alors $p \circ k = h \circ (i \otimes J)$.
\emph{Suffisance} : on modifie d'abord $h$, par extension d'homotopie le long
de $i$, en une homotopie $\tilde{h}$ avec $\tilde{h} \circ (i \otimes J) = p
\circ k$ ; on relève ensuite, $p$ étant une fibration, la flèche $X' \otimes
\{0\} \cup X \otimes J \to Y$ définie par $\varphi$ et $k$ au-dessus de
$\tilde{h}$ ; l'extrémité du relèvement est $\overline{\varphi}$.
\emph{Exemple où la condition manque} : en topologie, $p : Y = P S^{1} \to
S^{1}$ la fibration des chemins issus du point base $*$, $i : \{0, 1\} \to
[0, 1]$, $g$ constante sur le chemin constant $c_{*}$, $g'$ le lacet de degré
$1$, $\varphi = c_{*}$ constante. Les deux triangles commutent à homotopie
près, $[0, 1]$ étant contractile ; mais un $\overline{\varphi}$ serait un
relèvement du lacet $g'$ joignant $c_{*}$ à $c_{*}$, alors que tout relèvement
issu de $c_{*}$ aboutit dans la composante des lacets de degré $1$ de la
fibre $\Omega S^{1}$. Ces vérifications et l'exemple sont les nôtres.}

\pagerange{21}{23}
\section*{Relèvements et le carré relatif (pages 21 à 23)}

\emph{Page 21.} Trois carrés au crayon, sans texte, sont les trois formes
du relèvement dans une catégorie de modèles : le relèvement des homotopies,
$A \to A \times I$ contre une fibration $p : X \to Y$ ; sa forme duale, avec
un objet de chemins $A^{I} \xrightarrow{\partial_{0}} A$ et une cofibration
$Y \to X$ ; et le carré de relèvement général $A \to X$, $B \to Y$.\footnote{Il
marque la flèche $A \to A \times I$ « cof. » : c'est une cofibration
\emph{triviale}, inclusion d'une extrémité du cylindre, et c'est pourquoi le
relèvement existe. Le second carré porte « fib. » à gauche de $A$, et sa
flèche $Y \to A^{I}$ est tracée en trait double. À droite du premier carré,
deux petits croquis, une flèche prise dans une lentille, figurent
vraisemblablement une homotopie entre deux flèches. Deux signes isolés,
dont $Y_{\times}$ lu avec doute, ne sont pas interprétés.}

\emph{Page 22.} Le carré de la page 20, avec $g_{0} = \varphi i$, $g_{1} = g$,
$g'_{0} = p \varphi$, $g'_{1} = g'$, et $\theta_{0} = \varphi$, $\theta_{1} =
\overline{\varphi}$ les deux diagonales. On dispose, par hypothèse,
d'homotopies simpliciales
\[
k : \varphi i \longrightarrow g, \qquad h : p \varphi \longrightarrow g',
\]
et tout revient à pouvoir les choisir \emph{compatibles} à $i$ et à $p$,
c'est-à-dire avec $p \circ k = h \circ (i \otimes J)$ ; si $Y'$ est l'objet
final, c'est automatique.\footnote{« by assumption » et « simplicial » sont
ajoutés au-dessus de la ligne ; la lettre de la première homotopie est tracée
comme un $h$, mais les diagrammes de cette page et de la page 23 la nomment
$k$. La page 23 porte deux esquisses au crayon du même diagramme, sur un
listing daté 82155, début juin 1982 ; la première a des étiquettes
différentes, en partie douteuses, la seconde est celle de la page 22. Une
première tentative, par une somme $(X \sqcup X) \sqcup (X' \sqcup X')$ vers
$X \otimes J \sqcup X' \otimes J$ marquée « cofib? », est barrée.} Le carré

\begin{tikzcd}
X \otimes J \arrow[r, "k"] \arrow[d, "i \otimes J"'] & Y \arrow[d, "p"] \\
X' \otimes J \arrow[r, "h"'] & Y'
\end{tikzcd}

n'est pas commutatif a priori, mais seulement sur les extrémités $X \sqcup X
\to X \otimes J$.

La moitié inférieure de la page, annulée par deux traits croisés, réduit le
problème à un relèvement : avec $u = (g_{0}, g_{1}) : X \sqcup X \to Y$, $j : X
\sqcup X \to X \otimes J$ l'inclusion des extrémités et $v = h \circ (i
\otimes J)$, chercher $k'$ avec $k' j = u$ et $p k' = v$. On connaît $k$ avec
$k j = u$, mais seulement $p k \overset{s}{\sim} v$ et non $p k = v$. Comme $j$
est une cofibration non triviale, rien ne garantit le relèvement ; il existe
exactement quand $p k$ et $v$ sont homotopes relativement à $X \sqcup X$, ce
qui est le critère de la page 20.\footnote{Sur la page, « trivial cofibration »
est écrit à gauche de $j$, précédé d'un « non » biffé, ou lu tel avec doute ;
$j$ est une cofibration et n'est pas triviale, l'inclusion des deux extrémités
du cylindre n'étant pas une équivalence. La lettre $s$ au-dessus de $\sim$
peut être celle de « simpliciale », comme les homotopies du début de la page ;
lecture nôtre. Le mot qui précède « $pk \neq v$ » est surchargé. La page se
termine par « (But in case $Y'$ is final object we are evidently
through!) ».}

\pagerange{24}{24}
\section*{L'injectivité de $(*)$ (page 24)}

« Revenant au problème initial » : pour une suite de cofibrations $X_{i}
\hookrightarrow X_{i+1}$ de limite $X_{\infty}$, l'application
$\mathrm{Hom}(X_{\infty}, Y) \to \varprojlim \mathrm{Hom}(X_{i}, Y)$ est
surjective, ce qu'établit la page 20. Pour l'injectivité, soient $f, f' :
X_{\infty} \to Y$ dont les restrictions $f_{i}$, $f'_{i}$ à $X_{i}$ sont
homotopes pour tout $i$ ; il faut montrer $f \sim f'$, c'est-à-dire étendre
$(f, f')$ de $A = X_{\infty} \sqcup X_{\infty}$ à $B = X_{\infty} \times
\Delta(1)$, sachant que c'est possible au-dessus de chaque $X_{i} \times
\Delta(1)$.\footnote{« surjection » et « extend » sont lus avec doute. Ici
$A$ et $B$ désignent des objets, non plus les catégories de la première
partie ; on garde ses lettres, l'usage étant local à la page.}

Avec $A_{i} = X_{i} \sqcup X_{i}$ et $B_{i} = X_{i} \times \Delta(1)$, on
procède pas à pas : étant donné une homotopie $h_{i}$ sur $B_{i}$, on la
prolonge à
\[
B'_{i+1} = B_{i} \sqcup_{A_{i}} A_{i+1} = \bigl(\Delta(1) \times X_{i}\bigr)
\cup \bigl(\{0, 1\} \times X_{i+1}\bigr) \subset B_{i+1},
\]
par $\varphi_{i+1} = (f_{i+1}, f'_{i+1})$ sur $A_{i+1}$, et l'on cherche
$h_{i+1}$ sur $B_{i+1}$ prolongeant ce $h'_{i+1}$.

\begin{tikzcd}[column sep=4em, row sep=3em]
A_{i+1} \arrow[r, hook] \arrow[dr, "\varphi_{i+1}"'] & B'_{i+1} \arrow[r, hook] \arrow[d, "h'_{i+1}"] & B_{i+1} \arrow[dl, dashed, "h_{i+1} ?"] \\
 & Y &
\end{tikzcd}

L'inclusion $B'_{i+1} \subset B_{i+1}$ est une cofibration qui n'est pas
triviale en général : une homotopie $f_{i+1} \sim f'_{i+1}$ existe, mais pas
forcément une qui prolonge $h_{i}$.\footnote{La page dessine aussi $A_{i}
\hookrightarrow B_{i}$, $A_{i} \hookrightarrow A_{i+1}$ et une flèche
$\varphi_{i}$ de $A_{i}$ vers $Y$, les lettres $\varphi_{i}$, $\varphi_{i+1}$
étant de lecture douteuse ; les inclusions sont tracées en traits à crochet
sans pointe. Les $\{0\}$ et $\{1\}$ de la réunion sont surchargés et lus avec
doute.} La page abstrait alors la question : soit $Z_{0} \hookrightarrow Z_{1}
\hookrightarrow Z_{2} \hookrightarrow \cdots$ et une flèche $Z_{0} \to Y$ qui
s'étend à chaque $Z_{i}$ ; s'étend-elle à $Z_{\infty}$ ? « Pas clair du tout »,
tout le point étant de voir que
\[
\varprojlim_{i} \pi_{0}\bigl(\text{extensions de } Z_{0} \to Y \text{ à }
Z_{i}\bigr) \neq \varnothing .
\]
« Demander à Ronnie Brown ! »\footnote{Le mot qui suit $\pi_{0}$ est souligné,
surchargé et illisible ; la paraphrase « extensions de $Z_{0} \to Y$ à
$Z_{i}$ », que l'indexation $(Z_{i} ; Y)$ appelle, est la nôtre. La flèche
de $Z_{0}$ vers $Y$ est tracée double, $Y \Leftarrow Z_{0}$. Ronald Brown, à
Bangor, était à cette époque l'un de ses correspondants. Le cas de
l'injectivité est $Z_{i} = A \cup B_{i}$ dans $B$.}

\emph{La réponse.} La réduction est juste, et la réponse est négative en
général. Les espaces d'extensions $\mathrm{Map}_{Z_{0}}(Z_{i}, Y)$ forment une
tour de fibrations ; une extension à $Z_{\infty}$ est un point de leur limite,
et l'application de $\pi_{0}$ de la limite vers la limite des $\pi_{0}$ est
surjective. Une extension existe donc si et seulement si la limite des
$\pi_{0}$ est non vide --- et une limite projective d'ensembles non vides peut
être vide. Pour $(*)$, cela donne la suite exacte de Milnor : pour une suite de
cofibrations et $Y$ fibrant, la surjection $(*)$ a pour défaut d'injectivité
un $\varprojlim^{1}$, celui des $\pi_{1}$ des espaces d'applications
$\mathrm{Map}(X_{i}, Y)$ pointés en les $f_{i}$. Ce $\varprojlim^{1}$ ne s'annule pas toujours. Soit
$X_{i} = S^{1}$ pour tout $i$, avec les applications de degré $d \geq 2$, rendues
cofibrations par télescope, et $Y = K(\mathbb{Z}, 2)$. Alors
$\mathrm{Hom}_{\mathrm{Hot}}(X_{i}, Y) = H^{2}(S^{1}; \mathbb{Z}) = 0$, tandis
que
\[
\mathrm{Hom}_{\mathrm{Hot}}(X_{\infty}, Y) = H^{2}(X_{\infty}; \mathbb{Z})
\simeq \varprojlim\nolimits^{1}\bigl(\cdots \xrightarrow{\;d\;} \mathbb{Z}
\xrightarrow{\;d\;} \mathbb{Z}\bigr) \simeq \widehat{\mathbb{Z}}_{d}/\mathbb{Z} \neq 0 ,
\]
où $\widehat{\mathbb{Z}}_{d}$ est le complété $d$-adique de $\mathbb{Z}$. Donc $(*)$ n'est
pas injective : $X_{\infty}$ n'est pas limite inductive des
$X_{i}$ dans $\mathrm{Hot}$, et $\widehat{\Delta} \to \mathrm{Hot}$ ne commute
pas aux limites inductives dénombrables filtrantes, même le long de
monomorphismes.\footnote{Ce paragraphe est entièrement le nôtre : le dossier
s'arrête sur la question. La suite exacte est de J. Milnor (« On axiomatic
homology theory », 1962) ; la surjectivité de $\pi_{0}$ de la limite d'une
tour de fibrations vers la limite des $\pi_{0}$ est chez Bousfield et Kan
(1972). Une flèche $X_{\infty} \to Y$ nulle sur chaque $X_{i}$ sans l'être
est l'analogue, pour une suite, de ce qu'on appelle une application fantôme.
L'exemple ne dit pas que $\mathrm{Hot}$ manque de limites inductives
dénombrables, seulement que $X_{\infty}$ n'en est pas une ; le dossier ne
pose pas cette question distincte.}

\end{document}
